\begin{center}
{\Large \bfseries \textcolor{lyred}{第四章 不定积分}}
\end{center}

\vspace{6pt}
{\large \bfseries \textcolor{lyred}{第4.1 节 不定积分的概念与性质}}

\vspace{4pt}
{\bfseries \textcolor{lyred}{一.原函数与不定积分的概念}}

\textcolor{lyred}{定义}.若在区间I 上,
()
()
F
x
f x
'
=
,即
()
()
dF x
f x dx
=
,则称
()
F x 为
()
f x ,或者 
()
f x dx 在I 上的\textcolor{lyred}{原函数}. 
\textcolor{lyred}{例}. (
)
ln
x
x
+
+
是
x
+
的原函数, (
)
ln
x
x
C
+
+
+
也是. 
\textcolor{lyred}{注}.(1)连续函数必有原函数. 
(2)若
()
F x 为
()
f x 的原函数,则对任意常数C ,
()
F x
C
+
也是. 
(3)若
()
F x 和
()
G x 均为
()
f x 的原函数,则必有
()
()
G x
F x
C
=
+
. 
\textcolor{lyred}{定义}.
()
f x 在区间I 上带有任意常数项的原函数称为它在I 上的\textcolor{lyred}{不定积分},记为 
()
f x dx
\int 
,即
()
()
()
()
f x dx
F x
C
F
x
f x
'
=
+
\Leftrightarrow 
=
\int 
. 
()
f x -\textcolor{lyred}{被积函数},
()
f x dx -\textcolor{lyred}{被积表达式}, x -\textcolor{lyred}{积分变量},\int -\textcolor{lyred}{积分号},C -\textcolor{lyred}{积分常数}; 
()
f x dx
\int 
的图形构成一个彼此不相交的平行曲线族,其中每条曲线
()
y
F x
C
=
+
均称为
()
f x 的\textcolor{lyred}{积分曲线}. 
\textcolor{lyred}{例}.
x
x dx
C
=
+
\int 
, cos
sin
xdx
x
C
=
+
\int 
,
ln
dx
x
C
x
=
+
\int 
. 
\vspace{4pt}
{\bfseries \textcolor{lyred}{二.不定积分与微分的关系}}

(1)
()
()
()
()
f x dx
F x
C
d F x
C
f x dx
=
+
\Leftrightarrow 
+
=




\int 
; 
(2)
()
()
d F x
C
F x
C
+
=
+




\int 
;
()
()
d
f x dx
f x dx

=


\int 
. 
\vspace{4pt}
{\bfseries \textcolor{lyred}{三.基本积分表}}

kdx
kx
C
=
+
\int 
,
(
)
x
x dx
C
\mu 
\mu 
\mu 
\mu 
+
=
+
\neq -
+
\int 
,
ln
dx
x
C
x =
+
\int 
,
arctan
dx
x
C
x =
+
+
\int 
, 
arcsin
dx
x
C
x
=
+
-
\int 
, cos
sin
xdx
x
C
=
+
\int 
, sin
cos
xdx
x
C
=-
+
\int 
, 
tan
cos
dx
x
C
x =
+
\int 
,
cot
sin
dx
x
C
x =-
+
\int 
, sec tan
sec
x
xdx
x
C
=
+
\int 
, 
csc cot
csc
x
xdx
x
C
=-
+
\int 
,
x
x
e dx
e
C
=
+
\int 
,
ln
x
x
a
a dx
C
a
=
+
\int 
, sh
ch
xdx
x
C
=
+
\int 
, 
ch
sh
xdx
x
C
=
+
\int 
. 
\textcolor{lyred}{注}.
xdx
x
C
=
+
\int 
,
dx
C
x
x
=-
+
\int 
,
dx
x
C
x
=
+
\int 
. 
\vspace{4pt}
{\bfseries \textcolor{lyred}{四.直接积分法(分项积分法)}}

\textcolor{lyred}{性质(线性性)}.
()
()
()
()
f
x
g x
dx
f
x dx
g x dx
\alpha 
\beta 
\alpha 
\beta 
\pm 
=
\pm 




\int 
\int 
\int 
. 
\textcolor{lyred}{例}. (
)
3ln
x
x
x
x
dx
dx
x
x
x
C
x
x
x
-
-
+
-
=
=
-
+
+
+
\int 
\int 
. 
\textcolor{lyred}{例}.
(
)
(
)
(
)
arctan
ln
x
x
x
x dx
dx
dx
x
x
C
x
x
x
x
x
x
+
+
+
+


=
=
+
=
+
+


+
+
+


\int 
\int 
\int 
. 
\textcolor{lyred}{例}.
x
x
x
x
x
dx
dx
x
dx
x
x
x
+
+
+
-
-+


=
=
-+
=


+
+
+


\int 
\int 
\int 
4arctan
x
x
x
C
-
+
+
. 
\textcolor{lyred}{例}.
1 1
arctan
x
x
x
dx
dx
x
dx
x
x
C
x
x
x
-+


=
=
-+
=
-
+
+


+
+
+


\int 
\int 
\int 
. 
\textcolor{lyred}{例}.
1 1
arctan
x
x
x
x
dx
dx
x
x
dx
x
x
C
x
x
x
+-


=
=
-
+-
=
-
+
-
+


+
+
+


\int 
\int 
\int 
. 
\textcolor{lyred}{例}.
(
)
(
)
arctan
dx
x
x dx
dx
x
C
x
x
x
x
x
x
x
+
-


=
=
-
=-
-
+


+
+
+


\int 
\int 
\int 
. 
\textcolor{lyred}{例}.
(
)
tan
sec
sec
tan
xdx
x
dx
xdx
dx
x
x
C
=
-
=
-
=
-
+
\int 
\int 
\int 
\int 
. 
\textcolor{lyred}{例}.
1 cos
sin
sin
x
x
dx
dx
x
x
C
-
=
=
-
+
\int 
\int 
. 
\textcolor{lyred}{例}.
sec
tan
cos2
2cos
dx
dx
xdx
x
C
x
x
=
=
=
+
+
\int 
\int 
\int 
. 
\textcolor{lyred}{例}.
(
)
sin
cos
sec
csc
tan
cot
sin
cos
sin
cos
dx
x
x dx
x
x dx
x
x
C
x
x
x
x
+
=
=
+
=
-
+
\int 
\int 
\int 
. 
\textcolor{lyred}{例}.
,
,
,
2,
x
x
x
x
x
e
C
x
e
C
x
e
dx
e
C
x
e
C
x
-
-
-


+
<
+
\le 
=
=


-
+
>
-
+
+
>


\int 
. 
\textcolor{lyred}{补充练习 }
1.设
()
x
f x dx
e
C
-
=
+
\int 
,求
(
)
ln
I
x f
x dx
=\int 
. 
解.
()(
)
x
x
f x
e
e
-
-
'
=
=-
,故
(
)
ln
x
I
x
e
dx
xdx
x
C
-
=
-
=-
=-
+
\int 
\int 
. 
2.设
(
)
ln
f
x
x
'
=+
,求
()
f x . 
解.令
ln
u
x
=
,于是
()
u
f
u
e
'
=+
,故
()
(
)
x
x
f x
e
dx
x
e
C
=
+
=
+
+
\int 
. 
3.设
f
x
x
'




=








,求
()
f x . 
解.
()
()
dx
f
x
f
x
f
u
f x
C
x
x
x
u
x
x
-
-





'
'
'
\cdots 
=
\Rightarrow 
=-
\Rightarrow 
=
\Rightarrow 
=-
=
+










\int 
. 
4.
(
)
(
)(
)
(
)
(
)
1 1
x x
dx
x
x
dx
x
dx
x
dx
+
=
+-
+
=
+
-
+
=
\int 
\int 
\int 
\int 
(
)
(
)
x
x
C
+
-
+
+
. 
5.
(
)
(
)
(
)
(
)
(
)
x
x
dx
dx
dx
dx
dx
x
x
x
x
x
+
-
=
=
-
+
=
+
+
+
+
+
\int 
\int 
\int 
\int 
\int 
(
)
ln
x
C
x
x
+
+
-
+
+
+
. 
\vspace{6pt}
{\large \bfseries \textcolor{lyred}{第4.2 节 换元积分法}}

\vspace{4pt}
{\bfseries \textcolor{lyred}{一.凑微分法(第一类换元法)}}

\textcolor{lyred}{例}.
(
)
(
)(
)
(
)(
)
cos 3
cos 3
1 3
cos 3
x
dx
x
x
dx
x
d
x
'
+
=
+
+
=
+
+
=
\int 
\int 
\int 
(
)
cos
sin
sin 3
u
x
udu
u
C
x
C
=
+=
+
=
+
+
\int 
. 
\textcolor{lyred}{例}.
(
)
(
)
1 ln 3
x
d
x
dx
dx
x
C
x
x
x
'
+
+
=
=
=
+
+
+
+
+
\int 
\int 
\int 
. 
\textcolor{lyred}{例}.
1 sec
sec
tan
1 cos
dx
x
x
x
x
dx
d
C
x =
=
=
+
+
\int 
\int 
\int 
. 
\textcolor{lyred}{例}.
(
)
(
)
d a
x
d a
x
dx
dx
a
x
a
a
x
a
x
a
a
x
a
a
x
+
-


=
+
=
-
=


-
+
-
+
-


\int 
\int 
\int 
\int 
ln
ln
ln
a
x
a
x
a
x
C
C
a
a
a
a
x
+
+
-
-
+
=
+
-
.(公式16) 
\textcolor{lyred}{例}.
1 arctan
dx
a
x
x
d
C
a
x
a
a
a
a
x
a
=
=
+
+


+



\int 
\int 
.(公式17) 
\textcolor{lyred}{注}.类似地,
arcsin
dx
x
C
a
a
x
=
+
-
\int 
.(公式18) 
\textcolor{lyred}{例}.
(
)
(
)
x
x dx
x d
x
x
C
-
=-
-
-
=-
-
+
\int 
\int 
. 
\textcolor{lyred}{例}.
(
)
(
)
(
)
1 1
x
x
dx
d x
x
x
C
x
x
+
-
=
+
=
+
-
+
+
+
+
\int 
\int 
. 
\textcolor{lyred}{例}.
(
)
(
)
(
)
1 ln 1
n
n
n
n
n
n
n
n
n
n
dx
dx
x
dx
dx
x
C
x
x
n
n
x
x
x
x
x
x
x
-
+=
=
=
=
+
+
+
+
+
+
\int 
\int 
\int 
\int 
. 
\textcolor{lyred}{例}.
cos
cos
sin
dx
d
C
x
x
x
x
x
=-
=-
+
\int 
\int 
. 
\textcolor{lyred}{例}.
2arcsin 2
dx
dx
d
x
x
C
x
x
x
x
x
=
=
=
+
-
-
-
\int 
\int 
\int 
;或者, 
(
)
(
)
arcsin
d x
dx
x
C
x
x
x
-
-
=
=
+
-
-
-
\int 
\int 
. 
\textcolor{lyred}{例}.
(
)
arctan
arctan
arctan
2 arctan
arctan
x
x
u
dx
d
x
du
ud
u
x
u
x
x
=
=
=
=
+
+
+
\int 
\int 
\int 
\int 
(
)
arctan
x
C
+
;或者,
(
)
(
)
arctan
arctan
arctan
x
x
dx
d x
x
C
x
x
x
=
=
+
+
+
\int 
\int 
. 
\textcolor{lyred}{例}.
(
)
3 1
3 1
3 1
2 1
x
x
u
u
x
xe
e
e
dx
d
x
du
e
d u
e
C
u
x
x
-
+
-
+
-
-
-
+
=
+
=
=
=-
+
+
+
\int 
\int 
\int 
\int 
; 
或者,
(
)
3 1
3 1
3 1
3 1
x
x
x
x
xe
e
dx
d
x
e
d
x
e
C
x
x
-
+
-
+
-
+
-
+
=
+
=
+
=-
+
+
+
\int 
\int 
\int 
. 
\textcolor{lyred}{例}.
(
)
(
)
2ln
1 ln 1
2ln
2ln
2ln
d
x
dx
x
C
x
x
x
+
=
=
+
+
+
+
\int 
\int 
. 
\textcolor{lyred}{例}.
x
x
x
e
x
e
x
e
e
dx
e de
e
C
+
=
=
+
\int 
\int 
. 
\textcolor{lyred}{例}.
arctan
x
x
x
x
x
dx
de
e
C
e
e
e
-=
=
+
+
+
\int 
\int 
. 
\textcolor{lyred}{例}.
(
)
(
)
ln
ln
x
n
nx
nx
n
nx
x
nx
n
dx
de
du
u
e
C
C
e
n
u
n
e
e
e
u
u
=
=
=
+
=
+
+
+
+
+
+
\int 
\int 
\int 
. 
\textcolor{lyred}{例}.
sin
cos
tan
ln cos
cos
cos
x
d
x
xdx
dx
x
C
x
x
=
=-
=-
+
\int 
\int 
\int 
.(公式19) 
类似地,
cos
sin
cot
ln sin
sin
sin
x
d
x
xdx
dx
x
C
x
x
=
=
=
+
\int 
\int 
\int 
.(公式20) 
\textcolor{lyred}{例}.
1 sin
cos
sec
tan
1 sin
cos
cos
cos
dx
x
d
x
dx
xdx
x
C
x
x
x
x
-
=
=
+
=
-
+
+
\int 
\int 
\int 
\int 
. 
\textcolor{lyred}{例}.
(
)
sin
sin
sin
1 cos
cos
cos
cos
xdx
x
xdx
x d
x
x
x
C
=
=-
-
=-
+
+
\int 
\int 
\int 
. 
\textcolor{lyred}{例}.
(
)
sin
cos
sin
cos
sin
sin
1 sin
sin
x
xdx
x
xd
x
x
x
d
x
=
=
-
=
\int 
\int 
\int 
(
)
sin
2sin
sin
sin
sin
sin
sin
x
x
x d
x
x
x
x
C
-
+
=
-
+
+
\int 
. 
\textcolor{lyred}{例}.
sin
1 sin
sec
ln
ln sec
tan
cos
1 sin
1 sin
dx
d
x
x
xdx
C
x
x
C
x
x
x
+
=
=
=
+
=
+
+
-
-
\int 
\int 
\int 
; 
或者,
(
)
sec
sec
tan
sec
ln sec
tan
sec
tan
x
x
x
xdx
dx
x
x
C
x
x
+
=
=
+
+
+
\int 
\int 
.(公式21) 
类似地, csc
ln csc
cot
xdx
x
x
C
=
-
+
\int 
.(公式22) 
\textcolor{lyred}{例}.
tan
sin
sin
sin
sin
sin
cos
sin
cos
sin
cos
sin
x
x
x
x
dx
dx
d
x
d
x
x
x
x
x
x
x
=
=
=
=
\int 
\int 
\int 
\int 
sin
ln
arctan
sin
sin
u
x
du
du
du
x
C
u
u
u
x
+
=
-
=
-
+
-
-
+
-
\int 
\int 
\int 
. 
\textcolor{lyred}{例}.
(
)
(
)
sin
cos
sin cos
sin
cos
1 cos
cos
dx
xdx
d
x
du
x
x
x
x
x
x
u
u
-
=
=
=-
=
-
-
\int 
\int 
\int 
\int 
ln
u
du
C
u
u
u
u
u
u
-


-
+
+
=
+
+
+


-
+


\int 
;或者, 
sin
cos
sin
sin
cos
sin cos
sin cos
cos
sin cos
dx
x
x
x
x
x
dx
dx
dx
x
x
x
x
x
x
x
+
+
=
=
+
=
\int 
\int 
\int 
\int 
sin
sin
ln csc
cot
cos
cos
sin
3cos
cos
x
x
dx
dx
dx
x
x
C
x
x
x
x
x
+
++
=
+
+
-
+
\int 
\int 
\int 
. 
\textcolor{lyred}{例}.
1 cos2
1 cos2
sin 4
sin 2
sin
cos
x
x
x
x
x
x
xdx
dx
C
-
+



=
=
-
+
+






\int 
\int 
. 
\textcolor{lyred}{例}.
cos
1 2sin
sin
cot
sin 2
sin
sin
x
x
x
x
dx
dx
x
x
x
C
x
x
-
+
=
=-
-
+
-
+
\int 
\int 
. 
\textcolor{lyred}{例}.
(
)
sec
sec
sec
tan
tan
cos
dx
xdx
x
xdx
x
d
x
x =
=
=
+
=
\int 
\int 
\int 
\int 
tan
tan
tan
x
x
x
C
+
+
+
. 
\textcolor{lyred}{例}.
(
)
(
)
(
)
tan
sec
sec
tan
tan
sin
cos
tan
tan
tan
x
dx
x
x
dx
d
x
d
x
x
x
x
x
x
+
=
=
=
=
\int 
\int 
\int 
\int 
2tan
tan
tan
x
x
C
x
-
+
+
+
;或者,
sin
cos
sin
cos
sin
cos
dx
x
x dx
x
x
x
x
+
=
=
\int 
\int 
(
)
tan
tan
cos
sin
cos
cos
sin
dx
dx
dx
dx
x d
x
x
x
x
x
x
+
=
+
+
+
\int 
\int 
\int 
\int 
\int 
. 
\textcolor{lyred}{例}.
(
)
2tan
sec
2tan
arctan
4sin
9cos
4tan
4tan
d
x
dx
xdx
x
C
x
x
x
x
=
=
=
+
+
+
+
\int 
\int 
\int 
. 
\textcolor{lyred}{例}.
sec
tan
tan
sin
cos
tan
tan
dx
x
x
u
dx
d
x
du
x
x
x
x
u
+
+
=
=
=
=
+
+
+
+
\int 
\int 
\int 
\int 
tan
tan
arctan
x
u
x
du
d u
C
u
u
u
u
u
+
-


=
-
=
+






+
-
+




\int 
\int 
,或者, 
(
)
sin
cos
1 2sin
cos
sin 2
sin
2cos
d
x
dx
dx
du
x
x
x
x
x
u
u
=
=
=
+
-
-
+
\int 
\int 
\int 
\int 
. 
\textcolor{lyred}{例}.
(
)
(
)
(
)
tan
sec
tan
sec
sec
sec
sec
sec
x
xdx
x
xd
x
x
xd
x
=
=
-
=
\int 
\int 
\int 
sec
sec
sec
x
x
x
C
-
+
+
. 
\vspace{4pt}
{\bfseries \textcolor{lyred}{二.变量替换法(第二类换元法)}}

\textcolor{lyred}{例}.求
I
a
x dx
=
-
\int 
. 
解.令
sin
x
a
t
=
,
,
2 2
t
\pi \pi 


\in -




,则
2 1 cos2
cos
sin
t
I
a
td
t
a
dt
+
=
=
=
\int 
\int 
sin 2
sin cos
arcsin
a
a
a
a
a
x
x
t
t
C
t
t
t
C
a
x
C
a
+
+
=
+
+
=
+
-
+
. 
\textcolor{lyred}{例}.求
(
)
dx
I
x
x
=
+
-
\int 
. 
解.令
sin
x
t
=
,
,
2 2
t
\pi \pi 


\in -




,则
(
)
sin
1 sin
1 sin
cos
d
t
dt
I
t
t
t
=
=
=
+
+
\int 
\int 
(
)
arctan
2 tan
sec
tan
arctan
sec
tan
2tan
t
tdt
d
t
x
C
C
t
t
t
x


=
=
+
=
+




+
+
-


\int 
\int 
. 
\textcolor{lyred}{例}.求
dx
I
x
a
=
+
\int 
.(公式23) 
解.令
tan
x
a
t
=
,
,
2 2
t
\pi \pi 


\in -




,则
(
)
tan
sec
sec
d a
t
I
tdt
a
t
=
=
=
\int 
\int 
ln sec
tan
ln
ln
x
a
x
t
t
C
C
x
x
a
C
a
a
+
+
+
=
+
+
=
+
+
+
. 
\textcolor{lyred}{注}.
(
)
(
)
1 ln 2
d
x
dx
x
x
x
C
x
x
x
-
=
=
-+
-
+
+
-
+
-
+
\int 
\int 
. 
\textcolor{lyred}{例}.求
dx
I
x
a
=
-
\int 
.(公式24) 
解.当x
a
>
时,令
sec
x
a
t
=
,
0, 2
t
\pi 


\in 



,则
(
)
sec
sec
tan
d a
t
I
tdt
a
t
=
=
=
\int 
\int 
ln sec
tan
ln
ln
x
x
a
t
t
C
C
x
x
a
C
a
a
-
+
+
=
+
+
=
+
-
+
; 
当x
a
<-时,
ln
ln
u
x
du
I
C
x
x
a
C
u
a
u
u
a
=-
-
=
=
+
=
+
-
+
-
+
-
\int 
. 
\textcolor{lyred}{例}.求
dx
I
x x
=
-
\int 
. 
解.令
sec
x
t
=
,则
sec
arccos
sec tan
d
t
I
dt
t
C
C
t
t
x
=
=
\cdots 
=+
=
+
\int 
\int 
; 
或者,
arctan
1 1
xdx
d
x
I
x
C
x
x
x
-
=
=
=
-+
-+
-
\int 
\int 
; 
或者,
arcsin
dx
I
d
C
x
x
x
x
x
=
=-
=-
+
-
-
\int 
\int 
. 
\textcolor{lyred}{例}.求
x
I
dx
x
-
=\int 
. 
解.令
x
t
=
,则
(
)
(
)
3 2
x
I
t t
dt
t
d t
C
x
-
=-
-
=-
-
-
=-
+
\int 
\int 
. 
\textcolor{lyred}{例}.求
dx
I
x
x
=
+
\int 
. 
解.令
x
t
=
,则
tdt
x
I
t
C
C
x
t
+
=-
=-
+
+
=-
+
+
\int 
. 
\textcolor{lyred}{补充练习 }
1.
(
)
(
)
(
)
x dx
x
d
x
x
x
x
C
x
=
+
-
+
=
+
-
+
+
+
+
+
\int 
\int 
. 
2.
(
)(
)
x
x
x
x
dx
x
dx
dx
dx
x
x
x
x
x
x
+
+
-
+
=
=
+
=
+
+
+
+
-
+
\int 
\int 
\int 
\int 
arctan
arctan
x
x
C
+
+
. 
3.
arcsin
x
dx
dx
x
C
x
x
=
=
+
-
-
\int 
\int 
. 
4.
(
)(
)
(
)
(
)
(
)
(
)
(
)
(
)
1 ln
x
d x
d x
xdx
x
x
x
x
x
x
+
-
+
+
=
=
+
+
-
=
+
+
+
+
+
+
\int 
\int 
\int 
(
)
ln
arctan
x
x
x
C
+
+
+
-
+
. 
5.
(
)
sin 2
2sin
2sin
1 cos
4sin cos
8sin
cos
dx
dx
dx
dx
x
x
x
x
x
x
x
x
=
=
=
=
+
+
\int 
\int 
\int 
\int 
sec
tan
tan
tan
ln tan
8tan
tan
x
x
x
x
x
dx
d
C
x
x




=
+
=
+
+








\int 
\int 
. 
6.
ln tan
ln tan
tan
ln tan
sin cos
tan
x
x
dx
d
x
x
C
x
x
x
=
=
+
\int 
\int 
. 
\vspace{6pt}
{\large \bfseries \textcolor{lyred}{第4.3 节 分部积分法}}

\textcolor{lyred}{公式}.设()
u x , ()
v x 连续可导,则
()
()
()()
()
()
u x dv x
u x v x
v x du x
=
-
\int 
\int 
,即 
()()
()()
()()
u x v x dx
u x v x
u
x v x dx
'
'
=
-
\int 
\int 
. 
\textcolor{lyred}{类型一}.
()sin
nP
x
xdx
\int 
. 
\textcolor{lyred}{例}. (
)
(
)
(
)
1 cos
sin
1 sin
sin
x
xdx
x
d
x
x
x
x
xdx
+
=
+
=
+
-
=
\int 
\int 
\int 
(
)
(
)
1 sin
cos
1 sin
2 cos
2sin
x
x
xd
x
x
x
x
x
x
C
+
+
=
+
+
-
+
\int 
. 
\textcolor{lyred}{例}.
(
)
(
)
tan
sec
tan
x
xdx
x
x
dx
xd
x
x
=
-
=
-
=
\int 
\int 
\int 
(
)
(
)
tan
tan
tan
ln cos
x
x
x
x
x
x dx
x
x
x
C
-
-
-
=
+
-
+
\int 
. 
\textcolor{lyred}{例}.
tan
tan
2ln cos
cos
2cos 2
x
x
x
x
x
dx
dx
xd
x
C
x
x
=
=
=
+
+
+
\int 
\int 
\int 
. 
\textcolor{lyred}{例}.
cos
cos
sin
sin
cos
cos
cos
x
x
x
xdx
x
xd
x
xd
x
x
x




=-
=
-
=
-
-








\int 
\int 
\int 
cos
cos
sin
cos
cos
sin
sin
x
x
x
x dx
x
x
x
x
C






-
=
-
-
-
+
+












\int 
. 
\textcolor{lyred}{注}.
(
)
sin
sin
x
xdx
xd
xdx
=
\int 
\int 
\int 
. 
\textcolor{lyred}{例}.
1 sin
sin
1 sin
1 sin
cos
cos
x
x
x
dx
x
dx
x
dx
x
x
x
x
-


=
\cdots 
=
-
=


+
-


\int 
\int 
\int 
(
)
tan
tan
sec
ln cos
ln sec
tan
cos
x d
x
x
x
x
x
x
x
C
x


\cdots 
-
=
-
+
+
+
+




\int 
. 
\textcolor{lyred}{注}.
1 sin
1 sin
x
dx
xd
dx
x
x


=


+
+


\int 
\int 
\int 
. 
\textcolor{lyred}{类型二}.
()
x
nP x e dx
\int 
. 
\textcolor{lyred}{例}.
x
x
x
x
x
x
x
x
x e dx
x de
x e
e dx
x e
xe dx
x e
x de
=
=
-
=
-
=
-
\cdots 
=
\int 
\int 
\int 
\int 
\int 
(
)
x
x
x
x
x e
xe
e dx
x
x
e
C
-
+
=
-
+
+
\int 
. 
\textcolor{lyred}{例}.
(
)
x
x
x
x
x
x
x
xe dx
x
d e
xd e
x e
e
dx
e
e
=
-
=
-=
--
-
=
-
-
\int 
\int 
\int 
\int 
4arctan
x
x
x
x e
e
e
C
--
-+
-+
. 
\textcolor{lyred}{注}.
(
)
x
x
u t
x
x
x
u e
e
u
t
e
dx
d e
du
dt
e
u
t
=
=
-
-
-
=
-
=
=
=
+
+
\int 
\int 
\int 
\int 
2arctan
2arctan
x
x
dt
u
u
C
e
e
C
t


-
=
-
+
=
--
-+


+


\int 
. 
\textcolor{lyred}{例}.
(
)
x
x
x
x
x
x
x
xe
xe
xe
e
dx
xe d
e dx
e
C
C
x
x
x
x
x
-
=
=-
+
=-
+
+
=
+
+
+
+
+
+
\int 
\int 
\int 
; 
或者,
(
)
(
)
(
)
1 1
x
x
x
x
x
x
e
xe
e
e
dx
dx
dx
e d
C
x
x
x
x
x
+-
=
=
+
=
+
+
+
+
+
+
\int 
\int 
\int 
\int 
. 
\textcolor{lyred}{类型三}.
()ln
nP
x
xdx
\int 
. 
\textcolor{lyred}{例}.
3 ln
ln
ln
ln
x
x
x
x
x
x
xdx
xd
x
dx
x
C
=
=
-
=
-
+
\int 
\int 
\int 
. 
\textcolor{lyred}{例}.
(
)
(
)
ln cos
ln cos
tan
tan
ln cos
tan
ln cos
cos
x dx
xd
x
x
x
xd
x
x
=
=
-
=
\int 
\int 
\int 
tan
ln cos
tan
tan
ln cos
tan
x
x
xdx
x
x
x
x
C
+
=
+
-
+
\int 
. 
\textcolor{lyred}{例}.
(
)
(
)
ln
ln
x
x
x
dx
x
x
x
dx
x
+
+
=
+
+
-
=
+
\int 
\int 
(
)
(
)
(
)
ln
ln
2 1
d
x
x
x
x
x
x
x
x
C
x
+
+
+
-
=
+
+
-
+
+
+
\int 
. 
\textcolor{lyred}{例}.
(
)
ln
ln
ln
ln
2ln
xdx
x
x
x d
x
x
x
x dx
=
-
\cdots 
=
-
\cdots 
=
\int 
\int 
\int 
(
)
ln
2 ln
ln
ln
2 ln
x
x
x
x
x d
x
x
x
x
x
x
C
-
+
\cdots 
=
-
+
+
\int 
. 
\textcolor{lyred}{类型四}.
()arcsin
nP
x
xdx
\int 
. 
\textcolor{lyred}{例}.
arctan
arctan
arctan
arctan
x
x
x
x
x dx
x d
x
d
x
\cdots 
=
\cdots 
=
-
=
\int 
\int 
\int 
arctan
arctan
arctan
2 1
x
x
x
x
x
dx
x
x
C
x
-
=
-
+
+
+
\int 
. 
\textcolor{lyred}{例}.
(
)
arcsin
arcsin
arcsin
x
x
x
dx
d
x
x d
x
x
x
=-
-
=-
\cdots 
-
=
-
-
\int 
\int 
\int 
arcsin
arcsin
dx
x
x
x
x
x
x
C
x
-
-
+
-
=-
-
+
+
-
\int 
. 
\textcolor{lyred}{例}.
(
)
(
)
arctan
arctan
arctan
arctan
x
x
dx
x
dx
dx
x
x
x
x
x
x


=
-
=
-


+
+


\int 
\int 
\int 
, 
而
arctan
arctan
arctan
x
x
dx
x d
dx
x
x
x
x
x
-
=
\cdots 
=-
+
\cdots 
=
+
\int 
\int 
\int 
(
)
arctan
arctan
1 ln
x
x
x
dx
C
x
x
x
x
x
-
+
=-
+
+
+
+
\int 
. 
\textcolor{lyred}{例}.
(
)
arctan
arctan
arctan
ln
dx
x
x
dx
x
x
C
x
x
x
x
-
=
-
\cdots 
=
+
+
+
+
\int 
\int 
. 
\textcolor{lyred}{例}. (
)
(
)
arccos
arccos
2 arccos
x
x
dx
x
x
x
dx
x
-
=
-
\cdots 
=
-
\int 
\int 
(
)
(
)
arccos
2 arccos
arccos
2 1
arccos
x
x
xd
x
x
x
x
x
x
C
-
-
=
-
-
-
+
\int 
. 
\textcolor{lyred}{例}.
arcsin
sin
sin
sin
sin
sin 2
u
xdx
ud
u
u
u
udu
u
u
u
C
=
=
-
=
-
+
+
=
\int 
\int 
\int 
(
)
arcsin
x
x
x
x
C


-
+
-
+




. 
\textcolor{lyred}{例}.
(
)
(
)
arcsin
tan
tan
sec
x dx
ud
u
u
u
u
du
x
=
=
-
-
=
+
\int 
\int 
\int 
tan
tan
x
x
x
x
C
x
x
x
x
-
+
+
+
+
+
+
. 
\textcolor{lyred}{注}.以上两题也可以直接分部,但是会比较繁琐,故先换元简化,又例如: 
arccos
cos
sin
cos
cos
sin
sin
x
x
t
t
t
dx
d
t
t
t dt
td
t
t
x


=
=-
\cdots 
=-
-


-


\int 
\int 
\int 
\int 
. 
\textcolor{lyred}{注}. sin
sin
t
x
xdx
t
tdt
=
=
\int 
\int 
,
(
)
(
)
cos
1 cos
t
x
x
dx
t
tdt
=
-
-
=
+
\int 
\int 
, 
t
x
x
t
e
dx
t e dt
=
+
+
=
\int 
\int 
,
(
)
ln
2 ln
t
x
x dx
t
dt
x
=
+
=
-
+
\int 
\int 
.(先换元,再分部) 
\textcolor{lyred}{类型五(回归法},\textcolor{lyred}{循环法)}.
sin
xe
xdx
\int 
. 
\textcolor{lyred}{例}.
sin
sin
sin
cos
sin
cos
x
x
x
x
x
x
e
xdx
xde
e
x
e
xdx
e
x
xde
=
=
-
=
-
=
\int 
\int 
\int 
\int 
sin
cos
sin
x
x
x
e
x
e
x
e
xdx
-
-\int 
,故
(
)
sin
sin
cos
x
x
e
xdx
e
x
x
C
=
-
+
\int 
. 
\textcolor{lyred}{例}.
(
)
(
)
ln
sin ln
sin
sin
sin ln
cosln
t
x
t
t
x dx
t de
e
t dt
x
x
x
C
=
=
\cdots 
=
\cdots 
=
-
+
\int 
\int 
\int 
. 
\textcolor{lyred}{例}.
x
dx
x dx
x
x
dx
x
x
x dx
x
x
-
=
-
+
=
-
-
-
+
=
-
-
\int 
\int 
\int 
\int 
arcsin
x
x
x
C
-
+
+
. 
\textcolor{lyred}{例}.
(
)
sec
sec
tan
sec
tan
tan
sec
n
n
n
n
xdx
xd
x
x
x
n
x
xdx
-
-
-
=
=
-
-
=
\int 
\int 
\int 
(
)(
)
sec
tan
sec
sec
n
n
n
x
x
n
x
x dx
-
-
-
-
-
\int 
,故得递推公式: 
sec
sec
tan
sec
n
n
n
n
xdx
x
x
xdx
n
n
-
-
-
=
+
-
-
\int 
\int 
. 
\textcolor{lyred}{例}.
(
)
(
)
(
)
(
)
(
)
n
n
n
n
n
dx
x
x dx
x
dx
n
n
x
a
x
a
x
a
x
a
x
a
+
=
+
=
+
-
+
+
+
+
+
\int 
\int 
\int 
(
)
(
)
(
)
(
)
n
n
n
n
dx
dx
x
n
dx
na
na
x
a
x
a
na
x
a
x
a
+
+
-
\Rightarrow 
=
+
+
+
+
+
\int 
\int 
\int 
. 
\textcolor{lyred}{补充练习 }
1.求
(
)
ln
x
I
xdx
x x
+
=
-
\int 
.(提示:
(
)
(
)
(
)
1 2
x
x
x
x
x
x x
x x
x
+
+-
+
=
=
+
-
-
-
) 
解.
(
)
(
)
ln
2ln
ln
ln
2ln
2 ln
x
x
x
x
x
dx
I
dx
dx
xd
x
x
x
x x
x
-
=
+
=
+
=
-
+
=
-
-
-
-
\int 
\int 
\int 
\int 
ln
2ln
2ln
x
x
x
C
x
x
-
-
+
+
-
. 
2.求
(
)
arctan
x
e
I
dx
x
=
+
\int 
.(提示:令
tan
t
x
=
) 
解.
(
)
arctan
tan
cos
sin
cos
sec
2 1
t
t
t
x
e
e
x
I
d
t
e
t dt
t
t
C
e
C
t
x
+
=
=
\cdots 
=
+
+
=
+
+
\int 
\int 
. 
3.
x
dx
x dx
x
x
dx
x
x
x dx
x
x
+
=
+
-
=
+
-
+
+
=
+
+
\int 
\int 
\int 
\int 
(
)
ln
x
x
x
x
C
+
+
+
+
+
. 
\vspace{6pt}
{\large \bfseries \textcolor{lyred}{第4.4 节 有理函数的积分}}

\vspace{4pt}
{\bfseries \textcolor{lyred}{一.有理函数的积分}}

称
()
()
n
m
P x
Q
x 为\textcolor{lyred}{有理函数(分式)},当n
m
<
时,称为\textcolor{lyred}{真分式},当n
m
\ge 
时,称为\textcolor{lyred}{假分式}. 
有理函数总可以写成一个多项式与一个真分式的和,例如: 
x
x
x
x
x
x
+
+
+
=
+
+
+,
x
x
x
x
x
x
x
x
x
+
+
+
-
-+
=
=
-+
+
+
+. 
对于真分式
()
()
P x
Q x ,若
()
()
()
Q x
Q x Q
x
=
,其中
()
Q
x 与
()
Q
x 没有公因式,则 
()
()
()
()
()
()
P x
P x
P x
Q x
Q x
Q
x
=
+
,其中
()
()
()
()
,
P x
P x
Q
x
Q
x 均为真分式,称为\textcolor{lyred}{部分分式}. 
\textcolor{lyred}{定理}.设
()
(
)
(
)(
)
(
)
Q x
x
a
x
b
x
px
q
x
rx
s
\gamma 
\mu 
\alpha 
\beta 
=
-
-
+
+
+
+
\Lambda 
\Lambda 
,则真分式 
()
()
(
)
(
)
B
P x
A
A
B
Q x
x
a
x
b
x
a
x
b
\beta 
\alpha 
\alpha 
\beta 
=
+
+
+
+
+
+
+
-
-
-
-
\Lambda 
\Lambda 
\Lambda 
(
)
(
)
M x
N
R x
S
M x
N
R x
S
x
px
q
x
rx
s
x
px
q
x
rx
s
\gamma 
\gamma 
\mu 
\mu 
\gamma 
\mu 
+
+
+
+
+
+
+
+
+
+
+
+
+
+
+
+
+
+
\Lambda 
\Lambda 
\Lambda 
. 
\textcolor{lyred}{例}.求
x
I
dx
x
x
+
=
-
+
\int 
. 
解.设
x
A
B
x
x
x
x
+
=
+
-
+
-
-
,则
A
B
A
A
B
B
+
=
=-


\Rightarrow 


-
-
=
=


,因此,得 
x
x
x
x
x
+
-
=
+
-
+
-
-
;或者, (
)(
)
(
)(
)
x
x
x
x
x
x
+
-
+
=
=
-
-
-
-
x
x
x
x
x


+
-
=
-


-
-
-
-
-


,故
I
dx
dx
x
x
=
-
=
-
-
\int 
\int 
6ln
5ln
x
x
C
-
-
-
+
. 
\textcolor{lyred}{例}.求
(
)
dx
I
x x
=
-
\int 
. 
解.
(
)
(
)
(
)
A
B
C
x
x
x
x
x x
x
x
=
+
+
=
-
+
-
-
-
-
-
;或者, 
(
)
(
)
(
)
(
)
(
)
x
x
x x
x
x
x x
x x
x
x
-
+
=
=-
+
=
-
+
-
-
-
-
-
-
,故 
(
)
ln
x
I
dx
dx
dx
C
x
x
x
x
x
=
-
+
=
-
+
-
-
-
-
\int 
\int 
\int 
. 
\textcolor{lyred}{例}.求
(
)(
)
x
I
dx
x
x
x
+
=
+
+
+
\int 
. 
解.
(
)(
)
(
)(
)(
)
(
)(
)
A x
x
Bx
C
x
x
A
Bx
C
x
x
x
x
x
x
x
x
x
+
+
+
+
+
+
+
=
+
=
=
+
+
+
+
+
+
+
+
+
(
)
(
)
(
)(
)
A
B x
A
B
C x
A
C
x
x
x
+
+
+
+
+
+
+
+
+
,故
A
B
A
A
B
C
B
A
C
C
+
=
=




+
+
=\Rightarrow 
=-




+
=
=


,于是 
x
I
dx
x
x
x


=
-


+
+
+


\int 
,而
x
x
dx
d x
x
x
x
+
-


=
+
=


+
+




+
+




\int 
\int 
(
)
ln
arctan
ln
udu
du
u
u
C
x
x
u
u


-
=
+
-
\cdots 
+
=
+
+
-




+
+
\int 
\int 
arctan
x
C
++
,故
(
)
ln 2
ln
arctan
x
I
x
x
x
C
+
=
+
-
+
+
+
+
. 
\vspace{4pt}
{\bfseries \textcolor{lyred}{二.三角有理式的积分}}

对于它的积分
(
)
sin ,cos
R
x
x dx
\int 
,可以通过换元:
(
)
tan 2
x
t
x
\pi 
\pi 
=
-
<
<
,以及 
cos
t
x
t
-
=+
,
sin
t
x
t
=+
,
dt
dx
t
=+
,转化成为有理函数的积分.\textcolor{lyred}{ }
\textcolor{lyred}{例}.
(
)
(
)
1 sin
ln
sin
1 cos
t
x
dx
dt
t dt
t
t
t
C
x
x
t
t
+
+




=
=
+
+
=
+
+
+
=




+




\int 
\int 
\int 
ln tan
tan
tan
x
x
x
C
+
+
+
. 
\textcolor{lyred}{例}.
(
)
(
)
sin
cos
sin cos
1 sin
cos
sin
cos
sin
cos
x
x
x
x dx
dx
x
x
x
x
x
x
+
-
=
=
-
-
+
+
\int 
\int 
sin
cos
dx
x
x
+
\int 
,而
(
)
sin
cos
dx
dt
dt
x
x
t
t
t
=
=-
=
+
+
-
-
-
\int 
\int 
\int 
tan
ln
ln
2 2
tan
x
t
C
C
x
t
-+
-+
+
=
+
--
--
. 
\textcolor{lyred}{例}.
tan
sin
tan
tan
1 sin
sec
tan
2tan
t
x
x
x
x
t
dt
dx
dx
dx
x
x
x
x
t
t
=
=
=
=
\cdots 
=
+
+
+
+
+
\int 
\int 
\int 
\int 
(
)
arctan
arctan
arctan
2 tan
dt
dt
t
t
C
x
x
C
t
t
-
=
-
+
=
-
+
+
+
\int 
\int 
. 
\textcolor{lyred}{例}.
(
)
(
)
sin
cos
cos
sin
sin
cos
cos
sin
cos
sin
cos
sin
cos
x
x
x
x
d
x
x
xdx
x
dx
x
x
x
x
x
x
+
+
-
+
=
=
+
=
+
+
+
\int 
\int 
\int 
ln sin
cos
x
x
x
C
+
+
+
. 
\textcolor{lyred}{例}.求
sin
8cos
2sin
3cos
x
x
I
dx
x
x
+
=
+
\int 
. 
解.设
(
)
(
)
sin
8cos
2sin
3cos
2cos
3sin
x
x
A
x
x
B
x
x
+
=
+
+
-
, 2
A
B
A
B
-
=


+
=

,解得 
A
B
=


=

,故
2cos
3sin
ln 2sin
3cos
2sin
3cos
x
x
I
dx
x
x
x
C
x
x
-


=
+
=
+
+
+


+


\int 
. 
\vspace{4pt}
{\bfseries \textcolor{lyred}{三.简单无理函数}}

对于它的积分
,
ax
b
R x
dx
cx
d


+




+


\int 
,可以通过换元:
ax
b
t
cx
d
+
=
+
,转化为有理函数的 
积分. 
\textcolor{lyred}{例}.
(
)
x
u
x
u
u
dx
d u
du
du
x
u
u
u
-=
-


=
+
=
=
-
=


+
+
+


\int 
\int 
\int 
\int 
2arctan
2arctan
u
u
C
x
x
C
-
+
=
--
-+
. 
\textcolor{lyred}{例}.
(
)
x
u
d u
dx
u du
u
du
u
u
u
x
+=
-


=
=
=
-+
=


+
+
+


+
+
\int 
\int 
\int 
\int 
(
)
(
)
(
)
3ln 1
3ln 1
u
u
u
C
x
x
x
C
-
+
+
+
=
+
-
+
+
+
+
+
. 
\textcolor{lyred}{例}.
(
)
(
)
(
)
x t
dx
dt
t dt
t
dt
dt
t
t
t
t
t
t
x
x
=


=
=
=
=
-
=


+
+
+
+


+
\int 
\int 
\int 
\int 
\int 
(
)
(
)
arctan
arctan
t
t
C
x
x
C
-
+
=
-
+
. 
\textcolor{lyred}{例}.
(
)
(
)
(
)
(
)
x t
d
t
dx
t
t
dt
dt
t
t
t
t
t
t
t
x
x
x
-=
-
-
-
=-
=
=
=
+
+
+
-
-
+
-
\int 
\int 
\int 
\int 
(
)
6 1
12 1
12ln
t
dt
x
x
x
C
t


-
-+
=-
-
+
-
-
-
+
+


+


\int 
. 
\textcolor{lyred}{例}.
(
)
(
)
(
)
x
t
x
x
t
t
dx
t
td
t
t
dt
dt
x
x
t
t
t
+=
+
-
=
-
=
-
\cdots 
=-
=
-
-
-
\int 
\int 
\int 
\int 
ln
ln
t
x
x
x
dt
t
C
C
t
t
x
x
x
-
+
+
-


-
+
=-
-
+
=-
-
+


-
+
+
+


\int 
. 
\textcolor{lyred}{例}.
(
)(
)
(
)
x
t
x
dx
x
x
dx
dt
C
x
x
x
x
x
+=
-
+
+
=
=-
=-
+
-
-
-
+
-
\int 
\int 
\int 
. 
\textcolor{lyred}{例}.
arcsin
x
x
dx
dx
x
x
C
x
x
-
-
=
=
+
-
+
+
-
\int 
\int 
.\textcolor{lyred}{ }
\textcolor{lyred}{例}.求
dx
I
x
x
=
+
+
+
\int 
. 
解.令
x
x
t
+
+
=,
t
t
x
t
-
+
=
,则
I
dt
t
t
t


=
-+
-
=




\int 
(
)
ln
ln
x
t
t
C
x
x
x
x x
C
t
t


-
-+
+
=
-
+
+
+
-
+
+




. 
\textcolor{lyred}{注}.
(
)
ln
2arctan
x
t
e
x
t
e
dx
td
t
dt
t
t
C
t
=
-
-
=
+
=
=
-
+
=
+
\int 
\int 
\int 
2arctan
x
x
e
e
C
--
-+
. 
\textcolor{lyred}{补充练习 }
1.
sec
2 tan
2 tan
arctan
cos
2sec
2tan
dx
xdx
d
x
x
C
x
x
x
=
=
=
+
+
+
+
\int 
\int 
\int 
. 
2.
tan 2
tan
arctan
cos
x
u
x
dx
du
du
C
u
x
u
u
u
=
=
\cdots 
=
=
+
-
+
+
+
++
\int 
\int 
\int 
. 